\begin{table}[t]
\centering
\begin{threeparttable}
\caption{Independent enrichment ablation ($n=30$; query-bootstrap 95\% confidence intervals).}
\label{tab:ablation}
\small
\begin{tabularx}{\textwidth}{l*{3}{>{\centering\arraybackslash}X}}
\toprule
Method & Hit@5 & MRR & nDCG@5 \\
\midrule
BM25 raw & \textbf{0.867 [0.733, 0.967]} & \textbf{0.726 [0.593, 0.846]} & \textbf{0.662 [0.528, 0.789]} \\
R1: source chunks & 0.733 [0.567, 0.900] & 0.557 [0.412, 0.702] & 0.526 [0.380, 0.673] \\
R2: R1 + chunk summaries & 0.733 [0.567, 0.867] & 0.554 [0.408, 0.703] & 0.524 [0.380, 0.669] \\
R3: R2 + generated questions & 0.567 [0.400, 0.733] & 0.384 [0.253, 0.522] & 0.357 [0.220, 0.502] \\
R4: R3 + table summaries & 0.567 [0.400, 0.733] & 0.384 [0.252, 0.524] & 0.357 [0.217, 0.500] \\
BM25 + R4 RRF & 0.733 [0.567, 0.867] & 0.606 [0.461, 0.747] & 0.508 [0.366, 0.652] \\
Full adaptive & 0.800 [0.667, 0.933] & 0.608 [0.464, 0.747] & 0.571 [0.430, 0.713] \\
\bottomrule
\end{tabularx}
\begin{tablenotes}[flushleft]\footnotesize
\item R1--R4 are cumulative flat dense indexes. Generated representations map back to source chunks before evaluation.
\end{tablenotes}
\end{threeparttable}
\end{table}
